\bookchapter{Kapitel 17: Smartphone und Gerätevernetzung}

\chapterintro{Video-App direkt spiegeln}

\booksection{I. Synergie mit dem Computer: Smartphone und Computer sind nicht länger zwei verschiedene Welten}

Spiegelung funktioniert nicht?

\begin{itemize}
  \item \textbf{In 90 \% der Fälle sind beide Geräte nicht im selben Wi-Fi. In weiteren 10 \% ist die Spiegelungsfunktion des Fernsehers nicht aktiviert – suche in den Fernsehereinstellungen nach „Multi-Screen-Interaktion" oder „AirPlay" und schalte es ein.}III. Smartwatch/Fitness-Armband: Der Gesundheitsmanager am Handgelenk
  \item \textbf{Ich habe meiner Mutter ein Fitness-Armband gekauft. Anfangs sagte sie: „Ich treibe doch keinen Sport, wozu brauche ich das?" Nach zwei Wochen stellte sie fest, dass sie nur fünfeinhalb Stunden pro Nacht schlief, und begann, ihren Schlafrythmus aktiv anzupassen. Das ist der praktischste Wert von Smartwatch/Armband – sie macht unsichtbare Gesundheitsdaten zu sichtbaren Zahlen.}Benachrichtigungsfilter
  \item \textbf{: Wenn das Smartphone in der Tasche ununterbrochen vibriert, hebe das Handgelenk und wirf einen Blick – ist es eine WeChat-Nachricht vom Chef oder eine Pinduoduo-Werbung? In einer Sekunde entschieden; Nachrichten, für die sich das Herausnehmen des Smartphones nicht lohnt, lässt du einfach liegen.}Gesundheitsdaten-Synchronisation
\end{itemize}

\booksection{II. Spiegelung auf den Fernseher: Die ganze Familie schaut zusammen}

Transit-Karte/Zugangskarte synchronisieren

\begin{itemize}
  \item \textbf{:Auf dem Smartphone gebundene Transit- und Zugangskarten werden bei den meisten Marken automatisch auf die Uhr synchronisiert. In der U-Bahn muss man das Smartphone nicht herausnehmen – das Handgelenk an das Drehkreuz halten, und man ist durch. Man bekommt wirklich das Gefühl „Die Zukunft ist da".}Smartphone-Findungs-Wunder
  \item \textbf{: Auf der Uhr gibt es generell eine „Smartphone finden"-Funktion – ein Tippe, und das Smartphone klingelt. Ich nutze diese Funktion viel häufiger als erwartet – das Smartphone wird irgendwo hingelegt und sofort vergessen.}IV. Smart Home: Das Smartphone ist der „Hauptschalter" des Zuhauses
  \item \textbf{Smart Home klingt kompliziert, aber im Grunde sind es nur drei Dinge: Fernsteuerung per Smartphone, Automatisierungsregeln festlegen und Sprachsteuerung.}Einige meiner eigenen praktischen Anwendungen zeigen dir, wie nützlich das ist:
  \item \textbf{Klimaanlage}: Auf dem Heimweg im Sommer die Klimaanlage per Smartphone vorzeitig einschalten, und wenn man reinkommt, ist es schon kühl. Für etwas über hundert Yuan eine Klimaanlage-Partner-Steckdose kaufen, und das funktioniert – ohne die Klimaanlage auszutauschen.
\end{itemize}

\booksection{III. Smartwatch/Fitness-Armband: Der Gesundheitsmanager am Handgelenk}

: Vor dem Einschlafen im Bett, zu faul zum Aufstehen und Lichtausmachen – auf dem Smartphone tippen oder rufen: „Hey Siri / Xiao Ai, Licht aus!“

\begin{itemize}
  \item \textbf{Saugroboter}: Vor dem Verlassen des Hauses auf dem Smartphone starten, und wenn man zurückkommt, ist der Boden sauber. Man kann auch eine tägliche Zeitplanung einstellen.
  \item \textbf{Kamera}: Wenn zu Hause ältere Menschen, Kinder oder Haustiere sind, kann man jederzeit auf dem Smartphone einen Blick auf die Situation zu Hause werfen.
  \item \textbf{Welches Ökosystem wählen? Ein Satz fasst es zusammen:}Mit iPhone bevorzugt Geräte mit Apple HomeKit-Unterstützung wählen (etwas teurer, aber besserer Datenschutz).
  \item \textbf{Viele Geräte, Preis-Leistungs-Verhältnis wichtig Xiaomi-Mijia-Ökosystem ist am vollständigsten: Lampen, Steckdosen, Saugroboter – alles ist da.}Mit Huawei-Smartphone bevorzugt HarmonyOS Connect, funktioniert am besten mit Huawei-Geräten.
\end{itemize}

\booksection{IV. Smart Home: Das Smartphone ist der „Hauptschalter“ des Zuhauses}

Der Wechsel von Kabel- zu Bluetooth-Kopfhörern war für mich die spürbarste Upgrade-Erfahrung beim Smartphone. Nie wieder ein Knäuel aus der Tasche ziehen, nie wieder Sorge, dass das Kabel am Türknauf hängen bleibt.

Koppeln ist einfach

\begin{itemize}
  \item \textbf{: Beim ersten Einschalten des Kopfhörers wird normalerweise automatisch der Pairing-Modus gestartet (Blinken). Das Smartphone öffnet Bluetooth, findet ihn, ein Tippe – und verbunden. Danach verbindet er sich jedes Mal automatisch, wenn man ihn herausnimmt – man muss nichts weiter tun.}Gleiche Marke bringt Bonus
  \item \textbf{:AirPods mit iPhone zeigen automatisch ein Fenster mit dem Akkustand und wechseln nahtlos zwischen mehreren Apple-Geräten. Huawei FreeBuds mit Huawei-Smartphone und Xiaomi-Kopfhörer mit Xiaomi-Smartphone bieten ähnliche Pop-ups und Optimierungen.}Bluetooth-Lautsprecher
  \item \textbf{: Bei Familienfeiern oder beim Camping draußen kann ein Bluetooth-Lautsprecher die Stimmung verdoppeln. Sobald verbunden, wird dein Smartphone zur DJ-Station.}Latenz beachten
  \item \textbf{:Beim Spielen mit Bluetooth-Kopfhörern kann es zu leichter Latenz kommen (der Schuss ist erst eine halbe Sekunde später zu hören). Wer latenzempfindlich ist, sollte entweder Kabelkopfhörer verwenden oder Bluetooth-Kopfhörer mit „Low-Latency-Modus" kaufen.}VI. Mit dem Auto-System: Der sichere Beifahrer beim Fahren
\end{itemize}

Beim Fahren aufs Smartphone zu schauen ist illegal und lebensgefährlich. Aber manchmal braucht man tatsächlich Navigation, Telefonieren oder Musik. Hier kommen CarPlay (Apple) und Android Auto (Android) ins Spiel.

\begin{itemize}
  \item Wie es funktioniert
  \item : Wenn dein Auto es unterstützt (die meisten neueren Autos tun das), das Smartphone mit einem Datenkabel an den USB-Anschluss des Autos anschließen, und der Bildschirm wechselt automatisch zur CarPlay- oder Android Auto-Oberfläche. Einige neue Autos unterstützen drahtlose Verbindungen – noch bequemer.
  \item Was es kann
\end{itemize}

\booksection{V. Bluetooth-Kopfhörer und Lautsprecher: Die Freiheit des Kabellosen}

Auch alte Autos können das

\begin{itemize}
  \item \textbf{: Wenn dein Auto zu alt ist und dies nicht unterstützt, kannst du für ein paar hundert Yuan einen externen CarPlay-Bildschirm kaufen, der auf das Armaturenbrett geklebt wird – das Erlebnis ist ähnlich.}VII. Sicherheit und Datenschutz bei der Vernetzung
  \item \textbf{Je mehr Geräte vernetzt sind, desto größer ist die Angriffsfläche für Datenschutzverletzungen. Einige Punkte sind zu beachten:}Vor dem Spiegeln in der Öffentlichkeit zuerst prüfen, ob die Benachrichtigungsleiste des Smartphones private Nachrichten enthält – nicht die ganze Familie Dinge sehen lassen, die sie nichts angehen.
  \item \textbf{Bevor du ein Gerät weitergibst oder verkaufst, entferne die Kopplung – gehe in die Smartphone-Einstellungen, suche das Gerät und wähle „Entfernen" oder „Pairing aufheben".}Bei Smart-Home-Kameras Markenware mit Verschlüsselung kaufen, keine No-Name-Produkte für ein paar Dutzend Yuan. Es gibt zu viele Nachrichten über gehackte Kameras, die live übertragen wurden.
  \item \textbf{Probier es aus: Spiegle das Smartphone-Display auf den Fernseher und richte ein Smart-Home-Szenario ein.}Stelle sicher, dass Smartphone und Fernseher im selben Wi-Fi sind.
\end{itemize}

\begin{tcbinfo}{UWB-Technologie – Praeziseres „Raeumliches Gefuehl" als Bluetooth}
Zusaetzlich zu Bluetooth verbirgt sich in neueren Smartphones ein UWB-Chip (Ultra-Wideband). Seine Faehigkeit ist\textbf{praecise Entfernungsmessung und raeumliche Positionierung}– er berechnet, wie weit ein anderes Geraet entfernt ist und in welcher Richtung, mit Zentimetergenauigkeit. Bluetooth sagt nur „in der Naehe“; UWB sagt „1,2 m schraeg links vorne”. Haeufige Anwendungen:

\begin{itemize}
  \item \textbf{AirTag/„Suche"-Praezision}: Beim Suchen eines AirTag zeigt das iPhone Pfeilnavigation wie „7 m entfernt, nach links" – das funktioniert dank UWB. Bluetooth kann nur sagen, ob es in Reichweite ist.
  \item \textbf{Digitale Autoschlüssel}: UWB-faehige iPhones und Android-Flaggschiffe koennen als Autoschluessel dienen: Beim Naehern automatisch entriegeln, beim Weggehen verriegeln; unterscheidet Fahrer- und Beifahrerposition; schuetzt vor Relay-Angriffen.
  \item \textbf{Indoor-Navigation}: In Einkaufszentren und Flughäfen ermoeglichen UWB-Beacons „ganggenaue" Navigation – viel genauer als GPS.
  \item \textbf{Unterstützte Modelle}: iPhone 11 und spaeter haben UWB-Chips (U1/U2); bei Android sind Xiaomi, Samsung, vivo und OPPO in einigen Flaggschiffen dabei.
\end{itemize}

\end{tcbinfo}

\begin{tcbinfo}{LE Audio und Auracast – Die naechste Generation der Bluetooth-Audio-Standards}
Das zwanzig Jahre alte klassische Bluetooth-Audio (A2DP) wird durch\textbf{LE Audio}(Low Energy Audio) abgeloest – der neue Audiostandard ab Bluetooth 5.2:

\begin{itemize}
  \item \textbf{Stromsparender, bessere Klangqualitaet}: Mit dem neuen LC3-Codec; bessere Qualitaet als der alte SBC-Standard bei geringerem Verbrauch – deutlich laengere Akkulaufzeit der Kopfhoerer.
  \item \textbf{Dual-Stream-Audio}: Linke und rechte Ohrhoerer uebertragen unabhaengig; Wechsel in den Einzelohr-Modus ohne Abbruch.
  \item \textbf{Auracast Broadcast-Audio}: Das Aushaengeschild von LE Audio. Ein Smartphone/Fernseher kann gleichzeitig an\textbf{beliebig viele}Kopfhoerer senden – am Flughafen, im Museum, im Fitnessstudio oder im Hoersaal koennen viele Personen mit eigenen Kopfhoerern dasselbe hoeren.
  \item \textbf{Aktueller Stand}: 2024-2026 unterstuetzen neue Flaggschiff-Kopfhoerer und -Smartphones allgemein LE Audio; alte Geraete benoetigen ein Hardware-Upgrade, kein Software-Unlock. Beim Kauf neuer Kopfhoerer auf die LE-Audio-/Auracast-Kennzeichnung achten.
\end{itemize}

\end{tcbinfo}

\booksection{VI. Mit dem Auto-System: Der sichere Beifahrer beim Fahren}

Auf dem Smartphone: iPhone im Kontrollzentrum auf „Bildschirmspiegelung“ tippen; bei Android das Schnelleinstellungen-Panel nach unten ziehen und „drahtlose Spiegelung” suchen.

\begin{itemize}
  \item \textbf{Wähle dein Lieblingsfoto und spiegle es auf den Fernseher, damit die Familie es sieht.}Wenn du Smart-Home-Geräte hast, versuche in der Smartphone-App einen „Zu-Hause-weg-Modus" zu erstellen: Ein Tipp – Licht aus, Klimaanlage aus, Saugroboter startet.
  \item \textbf{Kapitel 18: Datenbackup und -migration}Eine Nachbarin ließ ihr Smartphone ins Wasserbecken fallen – nach dem Herausnehmen ließ es sich nicht mehr einschalten. Was sie am meisten bedauerte, war nicht das Smartphone selbst, sondern die Fotos von drei Jahren ohne Backup – die gesamten Aufzeichnungen vom Zeitpunkt der Geburt des Kindes bis zu seinen ersten Schritten. Als ich das hörte, wurde mir ganz übel, denn das hätte vollständig vermieden werden können. In diesem Kapitel richten wir ernsthaft eine Versicherung für deine Daten ein.
  \item \textbf{I. Zuerst klären: Was ist auf deinem Smartphone unwiederbringlich verloren?}Die „Überlebensstatus" verschiedener Daten sind unterschiedlich – du musst sie unterscheiden:
\end{itemize}

\booksection{VII. Sicherheit und Datenschutz bei der Vernetzung}

: Fotos und Videos (ohne Cloud-Sync), Notizen, SMS, Anrufprotokolle, heruntergeladene Dateien.

\begin{itemize}
  \item Auf App-Servern gespeichert – nach Gerätewechsel erneut einloggen und sie sind da
  \item : WeChat-Freundesliste, Taobao-Bestellungen, Transaktionsprotokolle in der Bank-App. Diese Dinge liegen auf den Servern der jeweiligen Unternehmen, nicht auf deinem Smartphone.
  \item In einer Grauzone – besondere Schritte erforderlich
\end{itemize}

\begin{tcbpractice}{: WeChat-Chatverläufe (auch wenn der WeChat-Server sie speichert, kannst du sie nicht selbst vom Server herunterladen), Geräteanmeldedaten der Mobile-Bank, Schlüssel in Zwei-Faktor-Authentifizierungs-Apps (wie Google Authenticator).}
\begin{enumerate}
  \item Du solltest also vor allem die erste und dritte Kategorie schützen.
  \item II. Vollständiges Backup: Alles auf einmal erledigen
  \item Ein vollständiges Backup ist wie ein vollständiger „Schnappschuss" des Smartphones – aktuelle Systemeinstellungen, Desktop-Layout, App-Daten, Fotos, Kontakte – alles eingepackt.
  \item Methode 1: Lokales Backup auf dem Computer (am zuverlässigsten, empfohlen für alle mit Computer)
  \item iPhone
\end{enumerate}

\end{tcbpractice}

\begin{tcbinfo}{Fallbeispiel 1: Akkulaufzeit-Optimierung in der Praxis}
\begin{itemize}
  \item \textbf{Szenario}: Der Nutzer eines Xiaomi 14 hat schlechte Akkulaufzeit – nach einem halben Tag ist der Akku leer.
  \item \textbf{Schritte}: (1) Einstellungen Akku Verbrauchsstatistik ansehen (2) App mit ungewoehnlich hohem Hintergrundverbrauch identifizieren (3) 5G ausschalten, auf 4G wechseln (4) Hintergrundaktivitaet einschraenken (5) Energiesparmodus aktivieren (6) Beim Laden Huelle abnehmen fuer bessere Kuehlung.
  \item \textbf{Erwartetes Ergebnis}: Akkulaufzeit verlaengert sich um 30-50 \%, Batteriezustand ist auslesbar.
\end{itemize}

\end{tcbinfo}
